Se diseñaron e implementaron seis pruebas funcionales para validar el esquema de cifrado multicapa (AES-256-CBC + AES-256-CTR + Spike + ChaCha20) en el escenario de drone de reparto. Todas las pruebas se ejecutaron en Python 3.13 sobre sistema Linux (Ubuntu 24.04), con procesador Intel Core i7 y 16 GB RAM.

\subsection*{A. Prueba 1: Generación de CRP con PUF Neuromórfica}

\textbf{Objetivo:} Verificar que la PUF neuromórfica con crossbar memristivo, neuronas LIF y codificación spike genera pares Reto-Respuesta deterministas y con la entropía esperada.

\textbf{Procedimiento:} Se instanció una \texttt{NeuromorphicHybridPUF} con parámetros $n=32$, $k=4$, semilla $=42$. Se generaron 10 CRP evaluando la respuesta con ruido desactivado y con ruido ($\sigma=0.02$).

\begin{table}[H]
    \centering\footnotesize\setlength{\tabcolsep}{3pt}
    \caption{CRP generados por la PUF Neuromórfica}
    \label{tab:crp-neuro}
    \resizebox{\columnwidth}{!}{\begin{tabular}{|c|c|c|c|}
        \hline
        \textbf{CRP \#} & \textbf{Reto (8 bits)} & $\mathbf{R_{\text{limpia}}}$ & $\mathbf{R_{\text{ruidosa}}}$ \\
        \hline
        1 & $[-1, -1, 1, 1, 1, -1, -1, -1]$ & +1 & +1 \\
        2 & $[-1, -1, 1, -1, 1, -1, 1, -1]$ & +1 & -1 \\
        3 & $[1, -1, -1, -1, -1, 1, -1, -1]$ & +1 & +1 \\
        4 & $[-1, -1, -1, -1, 1, -1, 1, 1]$ & -1 & -1 \\
        5 & $[1, -1, 1, 1, 1, -1, 1, -1]$ & -1 & -1 \\
        6 & $[1, 1, 1, -1, 1, 1, -1, -1]$ & +1 & +1 \\
        7 & $[-1, -1, -1, -1, 1, 1, 1, 1]$ & -1 & +1 \\
        8 & $[-1, -1, 1, 1, 1, -1, 1, 1]$ & +1 & +1 \\
        9 & $[1, -1, -1, -1, -1, -1, 1, 1]$ & -1 & -1 \\
        10 & $[1, 1, -1, 1, 1, 1, 1, 1]$ & +1 & -1 \\
        \hline
    \end{tabular}}
\end{table}

\textbf{Resultado:} La generación de CRP fue exitosa en el 100\% de los casos. La tasa de bit flip inducida por ruido fue del 30\% (3 de 10 respuestas cambiaron), consistente con el nivel de ruido configurado. La distribución de respuestas (60\% positivas) se mantuvo similar a la PUF Arbiter clásica, validando que la arquitectura neuromórfica preserva las propiedades estadísticas fundamentales de una PUF.

\subsection*{B. Prueba 2: Evolución STDP del Crossbar}

\textbf{Objetivo:} Verificar que el mecanismo STDP modifica las conductancias del crossbar tras cada autenticación exitosa.

\textbf{Procedimiento:} Se realizaron 10 autenticaciones consecutivas con la misma PUF neuromórfica, registrando la conductancia media del crossbar después de cada actualización STDP.

\begin{table}[H]
    \centering\footnotesize\setlength{\tabcolsep}{4pt}
    \caption{Evolución de conductancias con STDP}
    \label{tab:stdp-evol}
    \begin{tabular}{|c|c|c|}
        \hline
        \textbf{Auth \#} & $\mathbf{\mu_G}$ & $\mathbf{\sigma_G}$ \\
        \hline
        0 (inicial) & 0.500 & 0.150 \\
        1 & 0.503 & 0.149 \\
        2 & 0.507 & 0.148 \\
        3 & 0.510 & 0.147 \\
        4 & 0.514 & 0.146 \\
        5 & 0.517 & 0.145 \\
        6 & 0.521 & 0.144 \\
        7 & 0.525 & 0.143 \\
        8 & 0.528 & 0.142 \\
        9 & 0.532 & 0.141 \\
        10 & 0.535 & 0.140 \\
        \hline
    \end{tabular}
\end{table}

\textbf{Resultado:} La conductancia media aumentó progresivamente de 0.500 a 0.535 (+7\%) en 10 autenticaciones, mientras que la desviación estándar disminuyó de 0.150 a 0.140 (-6.7\%), indicando que el STDP refuerza selectivamente las sinapsis activas. Este comportamiento es consistente con el aprendizaje Hebbiano: ``neurons that fire together, wire together''.

\subsection*{C. Prueba 3: Cifrado y Descifrado AES-CBC + AES-CTR + Spike + ChaCha20}

\textbf{Objetivo:} Validar la integridad del cifrado cuádruple capa con cifrados modernos.

\textbf{Procedimiento:} Se cifró un mensaje de telemetría de 68 bytes con la cascada completa (AES-256-CBC $\to$ AES-256-CTR $\to$ Spike Permutation $\to$ ChaCha20) y se descifró, verificando la integridad byte exacta.

\begin{table}[H]
    \centering\footnotesize\setlength{\tabcolsep}{3pt}
    \caption{Parámetros de cifrado multicapa}
    \label{tab:enc-params-neuro}
    \resizebox{\columnwidth}{!}{\begin{tabular}{|l|l|}
        \hline
        \textbf{Parámetro} & \textbf{Valor} \\
        \hline
        Texto original & ``COMANDO: ENTREGAR PAQUETE ...'' (68 B) \\
        \hline
        AES-256-CBC IV & 8a1b2c3d4e5f6071... (16 B aleatorio) \\
        \hline
        AES-256-CTR Nonce & 9b2c3d4e5f607182... (16 B aleatorio) \\
        \hline
        Patrón Spike & $[0,1,0,1,1,0,1,0,\dots]$ (derivado del reto PUF) \\
        \hline
        IV ChaCha20 & 66bc52a63db636ec87bba07b747c65d5 (16 B) \\
        \hline
        Texto descifrado & Idéntico al original \\
        \hline
    \end{tabular}}
\end{table}

\textbf{Resultado:} 100\% de éxito en cifrado/descifrado sobre 100 iteraciones. La capa de permutación spike añade 0.3 ms adicionales al proceso de cifrado (2.4 ms total), un overhead aceptable.

\subsection*{D. Prueba 4: Simulación de Enlace Dron-Base con Cifrado Completo}

\textbf{Objetivo:} Validar el protocolo completo de autenticación, cifrado AES-CBC + AES-CTR + Spike + ChaCha20, y transmisión con evolución STDP.

\textbf{Procedimiento:} Se simularon 10 ciclos de autenticación + transmisión de telemetría entre dron y base, con actualización STDP tras cada autenticación exitosa.

\begin{table}[H]
    \centering\footnotesize\setlength{\tabcolsep}{3pt}
    \caption{Telemetría transmitida con cifrado multicapa}
    \label{tab:telemetry-neuro}
    \resizebox{\columnwidth}{!}{\begin{tabular}{|l|c|c|c|}
        \hline
        \textbf{Campo} & \textbf{Original} & \textbf{Descifrado} & \textbf{Estado} \\
        \hline
        drone\_id & NEURO-DRONE-01 & NEURO-DRONE-01 & $\checkmark$ \\
        \hline
        gps\_lat & -16.4090 & -16.4090 & $\checkmark$ \\
        \hline
        gps\_lon & -71.5374 & -71.5374 & $\checkmark$ \\
        \hline
        altitude & 150.0 & 150.0 & $\checkmark$ \\
        \hline
        speed & 12.5 & 12.5 & $\checkmark$ \\
        \hline
        battery & 85.0 & 85.0 & $\checkmark$ \\
        \hline
        heading & 270.0 & 270.0 & $\checkmark$ \\
        \hline
        G media crossbar & 0.521 & 0.521 & $\checkmark$ \\
        \hline
    \end{tabular}}
\end{table}

\subsection*{E. Prueba 5: Análisis de Entropía y Difusión}

\textbf{Objetivo:} Cuantificar la entropía de las respuestas de la PUF neuromórfica y verificar el efecto avalancha en la derivación de claves.

\textbf{Procedimiento:} Se generaron 1,000 CRP y se analizó la distribución de respuestas. Para el efecto avalancha, se midió el cambio en la clave ChaCha20 al modificar un solo bit del reto.

\textbf{Resultado:} La distribución de respuestas fue 49.7\% positivas y 50.3\% negativas ($\chi^2 = 0.18$, $p = 0.67$), consistente con una PUF perfectamente balanceada. El efecto avalancha mostró que cambiar un solo bit del reto produce un cambio promedio del 49.8\% en los bits de la clave derivada, muy cercano al ideal del 50\%. Esto valida que SHA-256 combinado con la PUF neuromórfica produce claves criptográficamente robustas.

\subsection*{F. Prueba 6: Detección de Ataque de Suplantación}

\textbf{Objetivo:} Demostrar que un dron falso con diferente PUF neuromórfica es detectado.

\textbf{Procedimiento:} Se crearon dron real (semilla 111) y falso (semilla 222). Se evaluaron 10 retos idénticos en ambas PUF.

\begin{table}[H]
    \centering\footnotesize\setlength{\tabcolsep}{3pt}
    \caption{Detección con PUF neuromórfica en 10 ejecuciones}
    \label{tab:attack-neuro}
    \begin{tabular}{|c|c|c|c|}
        \hline
        \textbf{Ejecución} & $\mathbf{R_{real}}$ & $\mathbf{R_{falso}}$ & \textbf{¿Coinciden?} \\
        \hline
        1 & +1 & -1 & No \\
        2 & -1 & +1 & No \\
        3 & +1 & -1 & No \\
        4 & -1 & -1 & Sí \\
        5 & +1 & -1 & No \\
        6 & -1 & +1 & No \\
        7 & +1 & -1 & No \\
        8 & +1 & -1 & No \\
        9 & -1 & +1 & No \\
        10 & +1 & +1 & Sí \\
        \hline
    \end{tabular}
\end{table}

\textbf{Análisis:} Coincidencia del 20\% (2/10), inferior al 30\% observado en la PUF Arbiter clásica. La PUF neuromórfica presenta mayor dispersión debido a la estocasticidad introducida por la codificación spike.

\subsection*{G. Prueba 7: Escalabilidad y Rendimiento}

\textbf{Objetivo:} Evaluar el rendimiento del sistema bajo carga creciente, simulando múltiples drones autenticándose simultáneamente.

\textbf{Procedimiento:} Se simularon 10, 50 y 100 drones conectándose concurrentemente a la estación base, midiendo el tiempo total de autenticación y el throughput.

\begin{table}[H]
    \centering\footnotesize\setlength{\tabcolsep}{3pt}
    \caption{Rendimiento del sistema bajo carga}
    \label{tab:scalability}
    \begin{tabular}{|c|c|c|c|}
        \hline
        \textbf{N° Drones} & \textbf{Tiempo total} & \textbf{Throughput} & \textbf{Latencia media} \\
        \hline
        10 & 24 ms & 416 auth/s & 2.4 ms \\
        50 & 127 ms & 393 auth/s & 2.5 ms \\
        100 & 268 ms & 373 auth/s & 2.7 ms \\
        \hline
    \end{tabular}
\end{table}

\textbf{Resultado:} El throughput se mantiene por encima de 370 autenticaciones por segundo incluso con 100 drones concurrentes, demostrando que el sistema escala linealmente. La latencia media aumenta solo 0.3 ms al quintuplicar la carga, indicando que el cuello de botella principal es la simulación de la PUF neuromórfica (2.1 ms por evaluación), no la concurrencia.

\subsection*{H. Análisis de Costo Computacional y Energético}

\textbf{Objetivo:} Cuantificar el costo computacional de cada componente del cifrado multicapa y estimar el consumo energético en una implementación hardware hipotética.

\textbf{Procedimiento:} Se midió el tiempo de ejecución de cada etapa del cifrado por separado utilizando \texttt{time.perf\_counter()} con 1,000 iteraciones. El consumo energético se estimó asumiendo un microcontrolador ARM Cortex-M4 a 100 MHz con 250 $\mu$A/MHz.

\begin{table}[H]
    \centering\footnotesize\setlength{\tabcolsep}{3pt}
    \caption{Desglose computacional del cifrado multicapa}
    \label{tab:cost-breakdown}
    \resizebox{\columnwidth}{!}{\begin{tabular}{|l|c|c|}
        \hline
        \textbf{Componente} & \textbf{Tiempo ($\mu$s)} & \textbf{Energía ($\mu$J)} \\
        \hline
        Codificación spike (rate coding) & 12 & 0.30 \\
        \hline
        Evaluación crossbar (4×64) & 45 & 1.13 \\
        \hline
        Dinámica LIF (10 pasos) & 380 & 9.50 \\
        \hline
        Actualización STDP & 28 & 0.70 \\
        \hline
        AES-256-CBC (relleno PKCS7, 240 B) & 42 & 1.05 \\
        \hline
        AES-256-CTR (240 B) & 38 & 0.95 \\
        \hline
        Permutación spike (240 B) & 22 & 0.55 \\
        \hline
        ChaCha20 (240 B) & 95 & 2.38 \\
        \hline
        Derivación de claves (3×SHA-256) & 18 & 0.45 \\
        \hline
        Total cifrado & 680 & 17.01 \\
        \hline
    \end{tabular}}
\end{table}

\textbf{Análisis:} La simulación de la dinámica LIF (10 pasos temporales) domina el costo computacional con el 56\% del tiempo total. En una implementación hardware con neuronas LIF analógicas reales, este costo se reduciría a nanosegundos. El AES-256-CBC y AES-256-CTR combinados representan el 12\% del tiempo, mostrando la eficiencia del hardware AES disponible en la mayoría de microcontroladores modernos. La energía total estimada de 17.01 $\mu$J por cifrado es compatible con la operación continua en un dron con batería de 5,000 mAh a 3.7 V (67 Wh), permitiendo aproximadamente 14 millones de cifrados por carga completa.

\subsection*{I. Vuelo Simulado Completo con Dashboard}

\textbf{Objetivo:} Validar el sistema completo con visualización en tiempo real de las señales neuromórficas.

\textbf{Procedimiento:} Se ejecutó el servidor Flask con 5 puntos de telemetría, verificando la recepción de eventos SSE con datos de conductancia, niveles spike y estado de neuronas LIF.

\begin{table}[H]
    \centering\footnotesize\setlength{\tabcolsep}{3pt}
    \caption{Vuelo simulado con monitoreo neuromórfico}
    \label{tab:flight-neuro}
    \begin{tabular}{|c|c|c|c|}
        \hline
        \textbf{Paso} & \textbf{$\mu_G$} & \textbf{Nivel Spike} & \textbf{Cifrado} \\
        \hline
        1 & 0.503 & 0.32 & 221 B \\
        2 & 0.508 & 0.45 & 227 B \\
        3 & 0.514 & 0.28 & 224 B \\
        4 & 0.521 & 0.61 & 226 B \\
        5 & 0.528 & 0.37 & 223 B \\
        \hline
    \end{tabular}
\end{table}

\subsection*{J. Tabla Resumen de Resultados}

\begin{table}[H]
    \centering\footnotesize\setlength{\tabcolsep}{3pt}
    \caption{Resumen de métricas del cifrado multicapa}
    \label{tab:summary-neuro}
    \resizebox{\columnwidth}{!}{\begin{tabular}{|l|c|}
        \hline
        \textbf{Métrica} & \textbf{Valor} \\
        \hline
        Precisión cifrado/descifrado & 100\% (1,000 iteraciones) \\
        \hline
        Prob. detección (1 bit) & $50.1\%$ \\
        \hline
        Prob. detección (64 bits) & $1 - 5.4 \times 10^{-20}$ \\
        \hline
        Tiempo cifrado total & 2.4 ms \\
        \hline
        Tamaño cifrado típico & 240--256 B \\
        \hline
        Sobrecarga de metadatos & 48 B (3 × 16 B IV/nonce) \\
        \hline
        Capas de cifrado & 4 (AES-CBC + AES-CTR + Spike + ChaCha20) \\
        \hline
        Claves independientes por capa & Sí (3 claves de 256 bits con sal) \\
        \hline
        Evolución de clave PUF & Sí (STDP post-autenticación) \\
        \hline
        Estándares empleados & AES-256 (NIST), ChaCha20 (NIST) \\
        \hline
    \end{tabular}}
\end{table}

